Un dipol està format per una càrrega positiva $+q$ i una càrrega negativa $-q$, del mateix valor, separades per 1,00~cm. En la figura s'han representat les superfícies equipotencials amb la mateixa separació de potencials entre cada parell de línies consecutives. Sabem que en el punt $P$ el potencial és de +10~V.

\figura[0.45]{fig1.png}

\begin{apartat}{1}
Reproduïu la figura i indiqueu els valors de potencial elèctric de cada una de les superfícies equipotencials que hi apareixen. Representeu-hi també, de manera aproximada, les línies de camp elèctric d'aquesta regió de l'espai.
\end{apartat}

\begin{apartat}{1}
Calculeu el valor de les càrregues $+q$ i $-q$.
\end{apartat}

\blocetiquetat{Dada:}{$k = \dfrac{1}{4\pi\varepsilon_0} = 8,99\times 10^{9}\ \mathrm{N\,m^2\,C^{-2}}$.}
